\chapter{Energy in Ecosystems}\label{ch:g7:eco-energy}

Sunlight powers nearly all \omterm{def:g7:eco-energy:ecosystem}{ecosystems}. \omterm{def:g2:food-energy:energy}{Energy} flows through \omterm{def:g7:eco-energy:chain}{food chains}
and webs; only part of it reaches each next level. This chapter builds
the \omterm{def:g6:nutrition:def}{energy} story of \omterm{def:g6:ecology:def}{ecology}.

\section{Ecosystem basics}

\begin{definition}[Ecosystem]\label{def:g7:eco-energy:ecosystem}
An \emph{ecosystem}\index{ecosystem} is a \omterm{def:g6:ecology:levels}{community} of \omterm{def:g4:ecosystems:levels}{organisms}
interacting with one another and with their non-living \omterm{def:g4:traits:env}{environment}
(\omterm{def:g4:plants-need:needs}{light}, \omterm{def:g4:plants-need:needs}{water}, soil, \omterm{def:g4:plants-need:needs}{air}, temperature).
\end{definition}

\begin{definition}[Producers and consumers]\label{def:g7:eco-energy:roles}
\emph{Producers}\index{producer} (mainly green \omterm{def:g1:plants-animals:plant}{plants} and algae) make
organic \omterm{def:g2:food-energy:food}{food} by \omterm{def:g7:photo-resp:photo}{photosynthesis}. \emph{Consumers}\index{consumer} eat other
\omterm{def:g4:ecosystems:levels}{organisms}: primary consumers eat producers; secondary consumers eat primary
consumers, and so on. \emph{Decomposers}\index{decomposer} (\omterm{def:g4:microbes:def}{fungi},
\omterm{def:g4:microbes:def}{bacteria}) break down dead material and recycle \omterm{def:g4:plants-need:needs}{nutrients}.
\end{definition}

\section{Food chains and webs}

\begin{definition}[Food chain and web]\label{def:g7:eco-energy:chain}
A \emph{food chain}\index{food chain} shows one feeding pathway. A
\emph{food web}\index{food web} shows many linked chains in a \omterm{def:g6:ecology:levels}{community}.
Arrows point in the direction of \omterm{def:g6:nutrition:def}{energy} flow (toward the eater).
\end{definition}

\omimg{food-web.jpg}{0.52\linewidth}{A \omterm{def:g7:eco-energy:chain}{food web} links many feeding
relationships in a \omterm{def:g6:ecology:levels}{community}.}

\begin{example}[Simple chain]\label{ex:g7:eco-energy:chain}
Grass $\to$ grasshopper $\to$ frog $\to$ heron. Grass is the \omterm{def:g7:eco-energy:roles}{producer};
the heron is a top \omterm{def:g7:eco-energy:roles}{consumer} in this chain.
\end{example}

\begin{omfigure}
\begin{tikzpicture}[font=\footnotesize,
  n/.style={draw, rounded corners, align=center, minimum width=1.8cm,
    minimum height=0.65cm}]
  \node[n, fill=omMeth!20] (g) at (0,0) {grass};
  \node[n, fill=omDef!15] (h) at (2.8,0) {grasshopper};
  \node[n, fill=omProp!15] (f) at (5.6,0) {frog};
  \node[n, fill=omThm!15] (b) at (8.2,0) {heron};
  \draw[-Stealth, thick] (g) -- (h);
  \draw[-Stealth, thick] (h) -- (f);
  \draw[-Stealth, thick] (f) -- (b);
  \node at (4.1,-0.9) {energy flows this way $\longrightarrow$};
\end{tikzpicture}

{\small \omterm{def:g7:eco-energy:chain}{Food chain}: arrows show \omterm{def:g6:nutrition:def}{energy} transfer, not ``who hunts whom''
wording alone.}
\end{omfigure}

\section{Energy pyramids}

\begin{proposition}[Energy loss between levels]\label{prop:g7:eco-energy:loss}
Only a fraction of \omterm{def:g6:nutrition:def}{energy} at one \omterm{def:g7:eco-energy:trophic}{trophic level} becomes biomass available
to the next (a common rule of thumb is about $10\%$). The rest is lost as
heat from respiration, incomplete consumption and waste. \omterm{def:g7:eco-energy:chain}{Food chains} are
therefore short.
\end{proposition}
\admitted

\begin{definition}[Trophic level]\label{def:g7:eco-energy:trophic}
A \emph{trophic level}\index{trophic level} is a feeding step: \omterm{def:g7:eco-energy:roles}{producers}
at level~1, primary \omterm{def:g7:eco-energy:roles}{consumers} at level~2, and so on.
\end{definition}

\begin{omfigure}
\begin{tikzpicture}[font=\small]
  \fill[omMeth!30] (0,0) rectangle (8,1.0);
  \node at (4,0.5) {producers (most energy in biomass)};
  \fill[omDef!25] (1.2,1.15) rectangle (6.8,2.0);
  \node at (4,1.55) {primary consumers};
  \fill[omProp!25] (2.2,2.15) rectangle (5.8,2.9);
  \node at (4,2.5) {secondary consumers};
  \fill[omThm!25] (3.0,3.05) rectangle (5.0,3.7);
  \node at (4,3.35) {top consumers};
  \node[font=\footnotesize] at (4,-0.45)
    {pyramid of energy (schematic)};
\end{tikzpicture}

{\small \omterm{def:g2:food-energy:energy}{Energy} available to each higher level is less --- the pyramid
narrows.}
\end{omfigure}

\begin{example}[Numbers]\label{ex:g7:eco-energy:10pct}
If \omterm{def:g7:eco-energy:roles}{producers} store \qty{10000}{kJ} of usable \omterm{def:g6:nutrition:def}{energy} in a \omterm{def:g1:seasons:season}{season}, about
\qty{1000}{kJ} might reach primary \omterm{def:g7:eco-energy:roles}{consumers} and about \qty{100}{kJ}
secondary \omterm{def:g7:eco-energy:roles}{consumers} under a strict $10\%$ rule --- a rough teaching model,
not a law for every \omterm{def:g7:eco-energy:ecosystem}{ecosystem}.
\end{example}

\begin{method}[Building a food web from a list]\label{met:g7:eco-energy:web}
\begin{enumerate}
  \item Identify \omterm{def:g7:eco-energy:roles}{producers}.
  \item Link each \omterm{def:g7:eco-energy:roles}{consumer} to what it eats; draw arrows toward the eater.
  \item Include \omterm{def:g7:eco-energy:roles}{decomposers} if data allow.
  \item Check that every \omterm{def:g7:eco-energy:roles}{consumer} has at least one \omterm{def:g6:nutrition:def}{energy} source.
\end{enumerate}
\end{method}

\section{Exercises}

\begin{exercise}[$\star$]\label{exo:g7:eco-energy:1}
Define \omterm{def:g7:eco-energy:ecosystem}{ecosystem}. Name two living and two non-living parts of a \omterm{def:g2:habitats:kinds}{forest}
\omterm{def:g7:eco-energy:ecosystem}{ecosystem}.
\end{exercise}

\begin{exercise}[$\star$]\label{exo:g7:eco-energy:2}
Define \omterm{def:g7:eco-energy:roles}{producer}, \omterm{def:g7:eco-energy:roles}{consumer} and \omterm{def:g7:eco-energy:roles}{decomposer}; give one example of each.
\end{exercise}

\begin{exercise}[$\star$]\label{exo:g7:eco-energy:3}
Redraw with correct arrow directions: hawk eats snake eats mouse eats
\omterm{def:g6:reproduction:plant}{seeds}.
\end{exercise}

\begin{exercise}[$\star$]\label{exo:g7:eco-energy:4}
What does a \omterm{def:g7:eco-energy:trophic}{trophic level} mean? Assign levels in grass $\to$ rabbit $\to$
fox.
\end{exercise}

\begin{exercise}[$\star$]\label{exo:g7:eco-energy:5}
Why are \omterm{def:g7:eco-energy:chain}{food chains} usually short?
\end{exercise}

